% 译者新增；不改变原文公式与表格编号。
\clearpage
\section*{翻译说明}
\phantomsection\addcontentsline{toc}{section}{翻译说明}
本译稿依据用户提供的 \texttt{arXiv-1509.00519v4.tar.gz} 及对应的 arXiv v4 PDF（2016 年 11 月 7 日）。正文、五个附录、原始图件、四张表和参考文献均予保留，采用翻译技能的蓝色标题中文单栏模板重排。下列已核实问题直接修入正文；本次复核不表示已穷尽原文全部问题。

\subsection*{数学与记号修正}
\begin{enumerate}
\item 第 2 节式 (1)：原文对隐藏状态作求和，但随后明确采用连续高斯潜变量。译稿将边缘化改为对全部潜变量的积分，并补上 $\mathrm{d}\hid$；离散潜变量时才使用求和。
\item 第 2 节式 (5)：原式左端写作 $\nabla_\params\log\expect_q[p/q]$，实际是边缘对数似然的梯度，并非所讨论的变分下界梯度。改为 $\nabla_\params\expect_q[\log(p/q)]$，与式~\eqref{eqn:vae_bound} 及后续重参数化推导一致。正文同时明确求导与期望交换的通常正则条件，以及重要性采样所需的支撑集覆盖条件。
\item 第 3.1 节式 (11)--(13)、(15) 及相关文字：源稿交替使用 $\hid(\aux,\obs,\params)$ 与 $\hid(\obs,\aux,\params)$。统一为第 2 节定义的 $\hid(\aux,\obs,\params)$；归一化权重分母的求和指标改用 $j$，避免与自由指标 $i$ 混用。
\item 定理 1 和附录 A 第三步：仅假设权重有上界，不能从强大数定律直接推出对数的期望收敛。例如，独立权重以相同概率取 $0$ 和 $2$，则 $\expect W=1$，但任何有限 $k$ 下平均权重为零的概率都为正，故 $\expect\log M_k=-\infty$。译稿采用明确的充分条件 $0<c\leq W\leq C<\infty$，并用 $|\log M_k|\leq\max\{|\log c|,|\log C|\}$ 与有界收敛定理补全证明。此条件是为了给出正确且简洁的陈述，并非收敛所必需的最弱条件；下界及单调性结论不需这一双侧界。
\item 第 3 节与附录 B：原文由取对数及 MAD 上界推断估计量不会有大方差，这一推断不成立。译稿保留正确的 $\mathrm{MAD}\leq2+2\delta$ 推导，补充 $\expect|\log\hat Z|<\infty$，并明确它不控制方差。例如，取非负随机变量 $Y$ 满足 $\expect Y<\infty$、$\mathrm{Var}(Y)=\infty$，令 $\hat Z=e^{-Y}/\expect[e^{-Y}]$，则 $\expect\hat Z=1$，$\log\hat Z$ 的 MAD 有限而方差无穷。
\end{enumerate}

\subsection*{实验文字与排版修正}
第 5.2 节称切换目标实验采用“表现最好的”模型，但表 2 的起始数值对应表 1 的单层 MNIST VAE（$k=1$，NLL 为 86.76）和 IWAE（$k=50$，NLL 为 84.78），并非表 1 全部模型中的最佳结果。译稿据表格明确所指模型，表中数据不改。

附录 D 表 3 的图注称 VAE 随样本数增加“略有受益”，但单层 VAE 的 NLL 为 88.71、88.83、89.05，且双层结果也非单调改善。译稿改为“未表现出一致改善”，保留全部原数值。

数学粗体使用 \verb|\bm|，微分 $\mathrm{d}$ 用正体；迹算子如出现，使用小写正体 $\operatorname{tr}$。原图完整保留，其中嵌入的英文标签不作覆盖重绘；可编辑网络图的文字标签已汉化。参考文献保留原文书目信息及作者年引用关系。拼写错误与表格多余空列在翻译排版时更正，不改变内容。
